\summarychapter{47}{Shift to the Global South}{What Is the ``New Christianity''?}

\begin{inbrief}
During the twentieth century most of the world's Christians came to live in Africa, Asia and Latin
America instead of Europe and North America. The chapter draws a profile of this southern
Christianity. It is Pentecostal or charismatic in practice; it inherited from Western missionaries a
voluntary and personal kind of faith; it reads the Bible as a book about its own world; it connects
the gospel with justice and with Christ's triumph over hostile powers; and it worships with the whole
body and sends out missionaries of its own.
\question{What kind of Christianity is growing in the Global South?}
\end{inbrief}

A map of the Christian world as D.~L.~Moody or Lenin knew it would be useless today. In 1900 nine
Christians in ten lived in Europe or North America; by 2000 at least seven in ten lived elsewhere.
The chapter gives examples. Nigeria's Anglican churches hold more worshipers in a week than all the
Anglican and Episcopal churches of Britain, Europe and North America combined. The Assemblies of God
are ten times larger in Latin America than in the United States, their country of origin. Congo has
more Baptists than Britain. And on any Sunday communist China has more people in church than western
Europe has, or North America.

The historian Philip Jenkins of Baylor University argued in \emph{The Atlantic Monthly} that religion
was taking the place of ideology as the main driving force in world affairs, and that a Christian
revolution was under way outside the liberal West. This revolution was turning toward the
supernatural and toward the New Testament's own outlook, which Jenkins summed up as ``a vision of
Jesus as the embodiment of divine power, who overcomes the evil forces that inflict calamity and
sickness upon the human race.'' With the faith's center of gravity in the former Third World, the
Christian globe had become ``pear-shaped.'' About the year 2000 Latin America counted 480 million
Christians, Africa 360 million and Asia 313 million, compared with 260 million in North America. The
liberal church leadership of the North Atlantic felt the consequences first at the Lambeth Conference
of 1998, when Anglican bishops from the south used their majority on behalf of views long out of
favor in the West; in Jenkins's phrase, ``Queen Victoria's ex-empire, from southern Africa to
Singapore, struck back.'' Africa after colonial rule shows the steepest climb, from 10 million
Christians among 107 million people in 1900 (some 9 percent) to 360 million among 784 million (46
percent) a hundred years later, and the proportion should keep rising, since Christian countries in
Africa have very high birth rates. Forecasts for the first quarter of the twenty-first century
reached 2.6 billion Christians worldwide, far more than any other religion.

\ssection{Global Christianity as a Pentecostal or Charismatic}
The first mark of the churches south of the equator is that the gifts of the Spirit have a regular
place in their services and in private prayer. To understand them one needs the history of
Pentecostalism. Its ancestors were Methodism and the Holiness movement, which had flourished on the
American frontier because they ignored distinctions of class, race and sex. Around 1900, after a
century and a half of Holiness teaching, Pentecostals took over that heritage and extended it into a
doctrine of three blessings: conversion; then sanctification, Wesley's theme; then baptism in the
Spirit, which equips the believer for service, as the Keswick teachers stressed, and which is shown
by speaking in tongues. Many believers linked this outpouring with the end of the age. Prophecy
conferences and renewed premillennial hopes raised expectations, healing ministries seemed to confirm
that God was making his people ready, and revivals in Wales, India and Korea lent support. Evan
Roberts, the voice of the Welsh revival, taught that no revival comes without such an experience of
the Spirit. The ``Korean Pentecost'' introduced customs that Korean churches have kept, including
prayer meetings at dawn and prayer spoken aloud by all at once.

\begin{sidebox}{Pentecostal and Charismatic}
People sometimes use the two words as synonyms. Strictly, \emph{Pentecostal} points to the Azusa
Street revival of the early twentieth century and to the denominations that came out of it, the
Assemblies of God and the Church of God for example, which generally teach that tongues are the sign
of receiving the Spirit. \emph{Charismatic} is the broader word. It is applied to anyone for whom the
more striking gifts (tongues, healing, prophecy) are a normal part of congregational worship and of
personal devotion, and it covers the teaching, practices and style of worship that began with
Pentecostalism and have since spread well outside those denominations.
\end{sidebox}

The usual account centers on two men. Charles Parham, who had been a Methodist minister, laid his
hands on a woman named Agnes Ozman, and she spoke in tongues; he took the language for Chinese.
Others had the same experience, Parham among them, whose tongue was supposed to be Swedish. He
believed that these were actual foreign languages given by miracle, the phenomenon called
\emph{xenolalia}, and that they were meant for missionary work overseas. William Seymour, a black
preacher, attended a three-month Bible course of Parham's at Houston, where segregation obliged him
to sit apart from the white students. He was then called to a black Holiness congregation in Los
Angeles, which would not accept what he taught about tongues. But at a meeting in the house of Edward
Lee, where he was staying, Seymour laid hands on Lee, who sank into a trance-like state and began to
speak in tongues, and seven others, Seymour included, received the same baptism that evening. Crowds
of every race soon overflowed the house, and the Azusa Street revival grew from those gatherings.

The book considers the rise and outreach of the Pentecostal and charismatic movements perhaps the
single most important development of the twentieth century for anyone who wants to understand the
church now. Pentecostal denominations are only a small part of it. The same impulse produced most of
today's nondenominational churches and has shaped many members of the older Protestant churches and
of the Catholic Church. American scholars count three waves: Azusa Street and the denominations it
generated; the charismatic renewal in mainline Protestantism and Catholicism during the 1960s and
early 1970s; and, from the 1980s, the acceptance of ``signs and wonders'' among conservative
Evangelicals, which started with John Wimber's teaching at Fuller Seminary in California and lives on
in the Vineyard churches. Historians have suggested wider settings as well. Azusa Street might be the
Jerusalem of a second Pentecost, or just one of several outpourings that happened in different
countries in the same years. The movement might be the largest part of a general turn to the Spirit
that also includes contemplative writers such as Henri Nouwen and Richard Foster. Or it might be one
episode in the long replacement of a fading liberalism by conservative and evangelical vigor. The
book allows each of these some truth and thinks the sheer size of the movement outweighs them.

Counting it is hard, because it runs through so many denominations. A cautious estimate is a little
over a quarter of all Christians, and the real figure may be much higher. The increase is the more
remarkable in a century when Christianity as a whole grew faster than ever. The book prints David
Barrett's regional figures for charismatic and Pentecostal Christians, here in millions:
\begin{sumtable}{L{0.44\linewidth}R{0.22\linewidth}R{0.22\linewidth}}
\sumhead{Region} & \sumhead{1900} & \sumhead{2000} \\ \midrule
Africa & 0.9 & 126 \\
Asia & 0 & 135 \\
Europe & 0 & 37.6 \\
Latin America & 0 & 141.4 \\
North America & 0 & 79.6 \\
Oceania & 0 & 4.3 \\
\end{sumtable}
It is only a small overstatement, scholars say, to call the worldwide growth of Christianity a growth
of Pentecostalism.

\ssection{Roots in Western Missionary Ventures}
Missionaries have often been charged with working for trade and empire as much as for Christ and
with confusing the gospel with their own customs. Jomo Kenyatta of Kenya gave the charge its classic
form: ``When the missionaries came to Africa they had the Bible and we had the land. They said `Let
us pray.' We closed our eyes. When we opened them we had the Bible and they had the land.'' Recent
research is kinder and shows that missionaries frequently pleaded for justice and for independence.
That they left a mark is beyond dispute. Mark Noll separates what they did directly from the pattern
they left, which lasted longer. The East African Revival serves as his test case. On the surface it
could not have been more Western. Joseph Church and his colleagues deliberately imitated the American
revivalist Charles Finney and read him aloud, and Church himself brought Keswick teaching, a Scofield
Bible, Western medicine and Western political ideas. It would still be a mistake to say that
foreigners engineered the revival. Many influences met in it, the fact that Africans listened to
these voices and not to others says something about Africa, and well-informed observers regard the
revival as truly African. In \emph{The New Shape of World Christianity} Noll therefore changes the
question. He asks less what the missionaries did than what form of Christianity they showed to the
people of the south.

That form was American. Christianity reached America as a European mission, but it changed there, and
the change explains the missionaries' deepest influence. Put briefly, American Christianity is
voluntary and owes less to Christendom.

\begin{sidebox}{Christendom}
At its simplest the word means the territory where Christians, as distinct from followers of other
faiths, live. It now usually has a cultural and political sense. A society may adopt Christian values
and write them into law; Sunday trading restrictions, the so-called blue laws, are an instance. For
some the word implies that Christianity created Western civilization. The general idea is that the
church as a religious arm and the government as a secular arm have separate tasks and yet cooperate
in a single order. Carried to its limit, in a state church, every Christian of the territory is a
citizen and every citizen is counted a Christian.
\end{sidebox}

No group rejected Christendom as completely as the Anabaptists. For them the church had no bond with
a state or a territory and consisted of those who freely joined it by covenant. The idea seems
obvious to Americans today. In the sixteenth century it was perilous, and Catholics and Protestants
alike hunted down those who held it and put them to death. Present-day Anabaptists maintain that
America has an unofficial Christendom of its own: no church is established, yet Christians think of
the country as somehow Christian and of the church as owing service to the nation, in the manner of a
chaplain. American believers thus stand somewhere between Anabaptist separation and the complete
Christendom of old Europe, a position the chapter considers essential for understanding them.

The kind of Christianity that spread most rapidly in North America, and that the Global South has
most readily adopted, was this voluntary kind. It also laid stress on the conversion of the single
person. Americans loved John Bunyan's \emph{Pilgrim's Progress} (1678), which remains Puritan and yet
tells the Christian life as one individual's conversion and journey, and the awakenings pressed the
same point. Southern churches have often accepted this emphasis although their own instincts are more
communal. American churches, further, were ready to follow leaders because of their gifts and
energy, whatever their rank; Wesley is remembered for his inventive organization and not for his
attempts to remain within the Church of England. What the missionaries passed on most durably,
the chapter concludes, was this voluntary, independent and inventive temper.

\ssection{Enacting the Bible}
Scripture speaks with unusual force in the south. Missionaries nearly always set about translating it
into the local tongue. The work gave them many chances to slip in their own assumptions, yet the
result, as several scholars have shown, was to strengthen independent thinking. People who receive
the Bible in their mother tongue gain confidence and a new respect for their own traditions,
especially when they meet similar customs in its pages. The chapter contrasts Islam, whose attachment
to the Arabic Koran works toward uniformity. Translation also let readers check their teachers.
Africans noticed that the patriarchs had several wives and asked why the missionaries were so strict
about monogamy.

Christianity became truly indigenous when readers made the Bible's speech, images and world their
own. Noll has described how differently this happens. Western Christians usually read the whole
through Paul's letters. Some Africans give that central place to Leviticus, because they recognize
its concern for holiness and ritual and learn from it that distinctions like clean and unclean are
biblical and need not be dismissed as heathen superstition. Some Asian believers turn first to
Proverbs and find there a God who values the wisdom they used to look for in Confucius. Southern
readers take the text as addressed to themselves, while Western readers are conscious of the
centuries that separate them from it. A Pentecostal does not find the New Testament world of tongues,
healings, demons and spiritual combat foreign. It is the pattern of normal Christian living.
Interpreters in the West tend to convert such things into terms they know and say, for example, that
a case once handled by exorcism would now be diagnosed as mental illness.

Conservative Westerners have their own ways of accounting for the distance. Classical
dispensationalism divided what the Bible says to the church from what it says to Israel, allowed no
mixture and regretted that Christians apply to themselves a Sermon on the Mount that was meant for
Israel. Cessationism goes further and teaches that miracles and charismatic gifts have come to an
end. Its best-known advocate, B.~B.~Warfield, a pillar of fundamentalist scholarship and an opponent
of the dispensationalists, held that miracles occurred in three clusters, each certifying a new stage
of revelation (the law given through Moses, the prophecy of Elijah and Elisha, and the kingdom
proclaimed by Jesus), and should not be expected now. The chapter remarks that such theories do for
very conservative Christians what liberal theories do for others: they relieve the unease of people
whose experience of church is unlike what they read in Scripture.

Behind this lies a difference about what is real. Believers in the Global South see the spiritual
and the material worlds constantly acting on each other. Westerners generally assume that the better
they control matter, through science for instance, the less they need spirit. The chapter imagines a
tribe whose shaman explains an epidemic as the gods' punishment for treachery toward an ally. A
Western scientist then analyzes the well, finds the water polluted and announces that the gods were
not involved. The reasoning rests on a false choice between natural and supernatural causes, which
the chapter compares to asking a child whether he will take the bus to school or take his lunch.
Western Christians often make concessions to this materialism. Afraid that knowledge leaves God
nothing to do, they let go of doctrines such as God's upholding of the world, or forget that God is
not simply one actor beside others; some abandon belief, and some keep a working deism in which God
remains in the wings and angels and demons are hardly mentioned. In the south the two realms meet,
so that spirits can affect how a person feels and fares, and readers there step into the biblical
text with fewer obstacles.

\ssection{The Gospel and Social Justice}
Most southern churches live among the poor, as the old term Third World suggests. Not all their
members are poor, but injustice is the setting they cannot escape, and they are more inclined than
northern Christians to believe that the gospel has to do with poverty. Liberation theologians teach
that Jesus aimed at deliverance from oppressive economic structures, find the Bible full of concern
for justice and demand that Christians do justice as well as preach it. Evangelicals in the south
listen to them more than North American Evangelicals do, without ceasing to put the saving work of
Jesus first.

Their situation even colors their view of that work. Christian theology has explained the atonement
in three principal ways. In the first, typical of Evangelicals, Christ stands in for sinners and
bears on the cross the penalty or debt that was theirs. In the second, favored by modern liberals,
the cross displays God's love so powerfully that it overcomes human distrust and calls forth love in
answer; what it cures is the inability to accept love and to give it. These two prevail in the West.
The third was the ordinary teaching of the Eastern church for a thousand years and has never died
out there. It sees humanity as enslaved: to an evil that includes a personal adversary, Satan, with
his demons; to sin, which no one overcomes unaided; and to death, which claims everyone. Left alone,
these powers would have the last word on every life. By coming, dying and rising, Jesus broke their
hold. His resurrection proves him the victor, and those united with him share his victory.

This view goes by the Latin name \emph{Christus Victor}, and southern Christianity is largely rooted
in it. Its prayers, its liturgies and its daily practice are full of triumph over enemies, and
believers there do not have to strain to imagine who the enemies are. Unstable governments,
corruption, violence and an economic oppression that seems to offer no exit are read as symptoms of a
damaged creation. In the Gospels Jesus attacks evil with power as a sign that the kingdom is coming,
and Paul writes of a creation that groans to be set free. The West has long neglected this
understanding of Christ's mission. In the south it underlies both theology and worship.

\ssection{Other Characteristics: Whole-Bodied Worship for the Whole World}
Non-Western churches have a feeling of solidarity that Westerners, who belong to the most
individualistic culture known to history, find hard to enter. A Western student grants that he is a
member of groups and yet thinks he has not explained himself until he has said what sets him apart.
Most of humanity does the reverse and identifies itself by naming the groups to which it belongs.
Entire communities may therefore convert together, and to be one body in the church is for them a
plain reality, where a Westerner hears a figure of speech. Worship shows it. The singing is lively,
uses the body and sounds noisy to Western ears, though it can go together with sacramental forms.
Intercession takes a large place, and special times of prayer may continue for hours. Many are
convinced that God speaks today as he spoke to the prophets. This does not push the Bible aside,
since a prophetic message often grows out of Bible study, from a phrase of the text. David Yonggi
Cho, who leads the largest congregation in the world, in South Korea, was asked why women share in
its leadership in so male-centered a society, and replied, ``God told me.''

These Christians also have a sharper awareness of evil and practice the spiritual warfare described
in Ephesians 6. Facing demonic manifestations during public worship is characteristic of world
Christianity generally and of Pentecostalism especially. Allan Anderson, a leading scholar of the
movement, traces its character and its growth to the belief that God has given healing back to the
church: ``in many cases, the ability of the preacher to heal is the primary cause of church
growth.'' Prayers for healing and the renunciation of evil spirits accordingly belong to the service.

Lastly, southern churches take up missionary work almost as soon as they exist. In common with
Pentecostals, many of them believe they are living at a decisive hour of the last days, when God is
gathering a people from every race to finish his work. The chapter's example is the house church
movement of China. Its first groups went on foot from one village to the next and were met with fury;
villagers turned out in a body to pelt the believers, whom they despised as cripples, beggars and
blind people, with stones and spoiled fruit. They persevered, lived underground for decades and
multiplied under Communist government. Many came to be known as the ``Back to Jerusalem'' movement, a
name Westerners often misread. It has nothing to do with evangelizing Jerusalem or Israel. The
thought is that the gospel set out from Jerusalem and traveled mostly westward, through North Africa
and Europe, and that the Great Commission will be complete when it has gone round the earth and come
back to its starting point. At the end of the twentieth century the movement hoped to send 100,000
missionaries into the fifty-one countries that lie between China and Jerusalem. Their route is the
ancient Silk Road, once the trade link between China and the Middle East, along which some 5,200
tribes and peoples, many of them Muslim, have not yet been reached.

\begin{keyfacts}
\keypeople{Philip Jenkins; Charles Parham; Agnes Ozman; William Seymour; Evan Roberts; John Wimber;
David Barrett; Mark Noll; Jomo Kenyatta; Joseph Church; John Bunyan; B.~B.~Warfield; David Yonggi Cho;
Allan Anderson.}
\keydates{%
  \kd{1678}{Bunyan's \emph{Pilgrim's Progress}}
  \kd{1900}{Ten percent of Christians live in the south and east; Africa about 9 percent Christian}
  \kd{1906}{Azusa revival: the ``first wave'' of Pentecostalism}
  \kd{1960s}{``Second wave'': charismatic renewal in mainline churches and Catholicism}
  \kd{1980s}{``Third wave'': signs and wonders among conservatives; John Wimber and the Vineyard}
  \kd{1998}{Lambeth Conference: southern Anglicans assert themselves}
  \kd{2000}{At least 70 percent of Christians live in the Global South; Africa 46 percent Christian}}
\keyterms{Global South; Pentecostal; charismatic; third blessing; Spirit baptism; \emph{xenolalia};
three waves; Christendom; voluntary Christianity; indigenization; cessationism; liberation theology;
\emph{Christus Victor}; spiritual warfare; ``Back to Jerusalem'' movement.}
\end{keyfacts}
